\chapter{Design: Intent-Aware Path Selection}
\label{ch:p1design}

Part~I asks how a single path should be chosen from a SCION host's candidate set so as to serve the
running application's intent under a changing network. This chapter presents the design. We start from
the DQN path selector of Bourak~et~al.~\cite{Bourak2025} and identify three limitations we set out to
remove (\cref{sec:p1design:strawman}). We then describe the improved simulation environment
(\cref{sec:p1design:env}), the state, action, and reward that define the decision problem
(\cref{sec:p1design:mdp}), the permutation-equivariant path-scoring architecture that handles a
variable path set (\cref{sec:p1design:scoring}), and finally the FiLM-based intent conditioning that is
the chapter's main contribution (\cref{sec:p1design:film}).

\section{Starting Point and Its Limitations}
\label{sec:p1design:strawman}

Bourak~et~al.~\cite{Bourak2025} frame SCION path selection as a reinforcement-learning problem and
train a DQN to pick a path from per-path measurements of latency, loss, and bandwidth, rewarding a
weighted combination of goodput and trust while charging for probing. They report oracle-level
selection quality at a fraction of the probing cost of measuring every path. This validates the RL
framing and gives us a reward structure to build on, but three design choices limit it.

\begin{description}
  \item[Fixed action space.] A vanilla DQN has one output per action, so the number of paths must be
  fixed and known. Real source--destination pairs have different, time-varying path counts; padding to
  a maximum wastes capacity and still caps the topology size.
  \item[Single baked-in objective.] The reward weights are constants chosen before training. The
  resulting policy embodies \emph{one} trade-off between goodput and trust; a different application
  intent demands a different, separately trained model.
  \item[Environment fidelity.] Capturing when path selection actually matters requires paths that
  genuinely share bottlenecks and conditions that vary over realistic timescales, so that a good choice
  at one hour is a poor choice at another.
\end{description}

Our design removes all three: a scoring architecture for the first, FiLM conditioning for the second,
and an improved environment for the third. A strawman for the second limitation -- simply appending the
reward weights to the state of the existing DQN -- is instructive precisely because it \emph{fails}, as
we show in \cref{sec:p1design:film}; that failure motivates the modulation-based design.

\section{Improved SCION Simulation Environment}
\label{sec:p1design:env}

Training and evaluating a path selector requires an environment in which the ``right'' path changes
with conditions. We construct one in several stages.

\paragraph{Topology.} AS-level topologies are generated with BRITE and converted into SCION structures
-- ISDs, core and non-core ASes, and inter/intra-ISD links -- so that the graph has realistic degree
and locality structure rather than being hand-drawn.

\paragraph{Control plane.} A beaconing simulation propagates path-segment construction beacons
top-down within each ISD and across the core mesh between ISDs, yielding the up/core/down segments that
compose into end-to-end paths. Peer links are excluded. The result is, for each source--destination
pair, a realistic set of candidate paths of varying length and quality -- the agent's action set.

\paragraph{Traffic and link dynamics.} We drive the network with 28~days of hourly demand: foreground
flows for every routable pair plus randomized background traffic, modulated by diurnal and weekly
factors. Crucially, load is aggregated \emph{per link}, so paths that share a bottleneck link contend
for its capacity. This is what makes selection non-trivial: the utilization and effective bandwidth of
a path depend on what every other flow is doing at that hour, and the best path at 3\,a.m. is often
congested at 8\,p.m.

\paragraph{Probing model.} Fresh metrics are not free. The environment charges a latency probe a small
per-hop cost and a full bandwidth probe a substantially larger one, and it credits the agent with
having probed only the path(s) it actually inspects. This turns ``measure less'' into a real objective
rather than an afterthought, and lets us compare the learned selector's probing footprint against
baselines that probe exhaustively.

\section{Decision Problem: State, Action, Reward}
\label{sec:p1design:mdp}

\paragraph{State.} The agent observes a \emph{global} context and a \emph{per-path} feature matrix. The
global context summarizes time and aggregate conditions -- normalized day-of-week and hour-of-day, mean
path utilization, mean trust, and the fraction of paths that are congested -- plus a compact encoding of
the source and destination. Each candidate path contributes a feature row: normalized latency, loss
rate, hop count, a bandwidth ratio, utilization, a static-bandwidth term, and a trust score. Because the
number of rows $N$ varies, the state is naturally a (global, path-matrix) pair rather than a flat
vector.

\paragraph{Action.} The action is the choice of one path, i.e.\ an index into the current candidate set.
Because $N$ varies, the action is realized as an $\arg\max$ over per-path scores (\cref{sec:p1design:scoring})
with invalid entries masked, rather than as a fixed categorical head.

\paragraph{Reward.} We adopt and generalize the goodput--trust reward of
Bourak~et~al.~\cite{Bourak2025}. For the chosen path, let $\mathrm{goodput}\in[0,1]$ be the achieved
bandwidth normalized by an auto-detected per-topology capacity, and define a trust term from loss and
delay,
\begin{equation}
\mathrm{trust} = \max\!\big(0,\; 1 - (w_3\,\mathrm{loss} + w_4\,\mathrm{delay})\big),
\label{eq:p1trust}
\end{equation}
with $\mathrm{delay}$ the latency capped at \SI{100}{\milli\second} and normalized. The base reward
trades goodput against trust and is shifted to $[-1,1]$, and a probing penalty is subtracted:
\begin{equation}
r = \underbrace{2\big(w_1\,\mathrm{goodput} + w_2\,\mathrm{trust}\big) - 1}_{\text{base reward}}
    \;-\; w_{\text{probe}}\cdot \mathrm{probe\_penalty},
\label{eq:p1reward}
\end{equation}
clamped to $[-1,1]$, where $\mathrm{probe\_penalty}$ is the probing cost normalized to $[0,1]$. The
five weights $\mathbf{w} = (w_1, w_2, w_3, w_4, w_{\text{probe}})$ are exactly the \emph{application
intent}: default balanced values are $(0.7, 0.3, 0.5, 0.5, 0.05)$, but a throughput-hungry intent
pushes $w_1$ toward~1, a loss-averse intent raises $w_2$ and $w_3$, a latency-sensitive intent raises
$w_4$, and a probe-frugal intent raises $w_{\text{probe}}$. The design goal of \cref{sec:p1design:film}
is a \emph{single} policy that behaves correctly for \emph{any} setting of $\mathbf{w}$.

\section{Handling a Variable Path Set: The Scoring Architecture}
\label{sec:p1design:scoring}

To make the policy independent of the path count, we score each path with a \emph{shared} network and
select by comparison, following the permutation-equivariant design of \cref{sec:background:variable}.
The global context is broadcast and concatenated to every path's feature row, and a shared multilayer
perceptron maps each $(\text{global}, \text{path}_i)$ pair to a scalar $Q$-value. Selection is
$\arg\max_i Q_i$; absent or invalid paths are masked to $-10^9$ so they can never be chosen and do not
distort training targets.

We use a \emph{dueling} scoring head (\cref{sec:background:rl}): a value stream reads the global context
alone to produce $V(s)$, an advantage stream reads the per-path features to produce $A(s, i)$, and the
two combine as $Q_i = V(s) + \big(A(s,i) - \operatorname{mean}_{j\in\text{valid}} A(s,j)\big)$ with the
mean taken only over valid paths. Because many paths are often comparably good, separating ``how good is
this situation'' from ``which path is marginally better'' stabilizes learning. Double DQN targets and
prioritized replay (\cref{sec:background:rl}) further stabilize and accelerate training. The scoring
design also delivers the dynamic-action property directly: the same weights serve a pair with three
paths and a pair with thirty, and are invariant to how the paths happen to be ordered.

\section{Intent Conditioning via Feature-Wise Modulation}
\label{sec:p1design:film}

The central contribution of Part~I is making one scoring policy serve every intent $\mathbf{w}$. We
condition the network on $\mathbf{w}$, but the \emph{mechanism} of conditioning is what makes or breaks
it.

\subsection{Why concatenation fails}
\label{sec:p1design:concat}
The obvious approach is to append $\mathbf{w}$ to the state. But consider a scoring policy: the intent
is part of the \emph{global} context, which is broadcast \emph{identically} to every path before
scoring. Adding the same constant vector to every path's input tends to move all path scores by a
similar amount -- a common-mode shift -- which leaves the $\arg\max$, and hence the chosen path,
essentially unchanged. Formally, if conditioning enters only additively and symmetrically across paths,
it cancels under the comparison that defines the action. The policy can thus appear to ``take intent as
input'' while its behavior is invariant to it. This is not a hypothetical; it is the failure mode our
ablation in \cref{ch:p1eval} is designed to expose, and it is why a naive goal-conditioned baseline is
inadequate.

\subsection{Weight-FiLM dueling head}
The fix is to let the intent modulate the \emph{per-path} features multiplicatively, so that it changes
the relative ranking of paths rather than a shared baseline. A small conditioning network maps the intent
vector to feature-wise scale and shift parameters,
\begin{equation}
(\gamma, \beta) = \mathrm{FiLM}_\theta(\mathbf{w}), \qquad
\tilde{f}_i = (1 + \gamma) \odot f_i + \beta,
\label{eq:p1film}
\end{equation}
where $f_i$ is path~$i$'s feature row and $\tilde{f}_i$ the modulated features fed to the advantage
stream. A latency-sensitive intent can learn a $\gamma$ that amplifies the latency feature, so paths
differ more sharply in their latency contribution and a low-latency path wins; a throughput-hungry intent
can instead amplify bandwidth. Because the modulation is per-feature and multiplicative, it acts
\emph{differently} on paths with different feature values and therefore genuinely reorders them -- exactly
what concatenation cannot do. The value stream continues to read only the topology/time context, keeping
the state-value estimate intent-agnostic while the advantage -- the part that decides \emph{which} path --
is intent-modulated.

\subsection{Training across intents}
To learn a policy that spans the intent space we train across a set of \emph{reward profiles} -- named
weight vectors spanning balanced, throughput-maximizing, trust/loss-averse, latency-sensitive, and
probe-frugal extremes. Each training episode samples a profile, sets the environment's reward weights to
it, and feeds the same weights to the FiLM head, so the agent simultaneously experiences the reward for an
intent and the encoding of that intent. Profiles are scheduled in a stratified fashion so every intent
receives balanced coverage. At inference, presenting a new $\mathbf{w}$ -- including interpolations between
trained profiles -- yields the corresponding behavior without any retraining.

\subsection{Relation to Mortise}
This is the path-selection instance of Mortise's principle~\cite{Shen2026}. Where Mortise maps an
application's revealed QoS preference to congestion-control \emph{parameters}, our FiLM head maps an
explicit intent vector to a \emph{modulation} of path features that determines which path is chosen. In
both, intent is a first-class input that reshapes a network mechanism, not a constant compiled into it.
\Cref{ch:p2design} pushes the same principle one step further, letting intent shape not only the network's
choice but the application's own output.
